\rtmtight
\begin{tblr}{width=\textwidth,colspec={|Q[c,m,2.1cm]|X[l,m]|},hlines,vlines,hspan=minimal,rowsep=1.5pt,colsep=3pt,row{1}={bg=headgray,font=\bfseries}}
\SetCell{c,m}\hfil\logo\hfil & \textbf{POLITEKNIK NEGERI MALANG}\par\textbf{JURUSAN TEKNOLOGI INFORMASI}\par\textbf{PROGRAM STUDI: D4 TEKNIK INFORMATIKA} \\
\SetCell[c=2]{l} \textbf{RENCANA TUGAS MAHASISWA (RTM) DAN RUBRIK PENILAIAN} & \\
Nama Mata Kuliah & Metodologi Penelitian \\
Kode Mata Kuliah & RTI255002 \quad \textbf{sks / Jam perminggu:} 2 SKS/4 jam \quad \textbf{Semester:} 5 \\
Dosen Pengampu & Tim Pengajar Program Studi D4 TI \\
\end{tblr}
\assessmentblock{Tugas Studi Literatur dan Perencanaan Penelitian}{Case Method dan penulisan ilmiah}{SCPMK0403-03801, SCPMK0403-03802, SCPMK0403-03803}{Mahasiswa menyelesaikan tugas-tugas terkait studi literatur, identifikasi masalah, dan perencanaan metode penelitian sesuai Sub-CPMK.}{Menganalisis kasus atau topik penelitian, menyusun artefak ilmiah, mendokumentasikan evidence, dan mempresentasikan hasilnya bila diminta.}{Laporan studi literatur, matriks sintesis, draft rumusan masalah, rencana metode, dan/atau bahan presentasi.}{Ketepatan identifikasi dan sintesis literatur ilmiah & 5 & Minimal 10 sumber relevan disintesis dengan tepat dan kerangka konseptual terdokumentasi secara sistematis. & Sumber cukup relevan tetapi sintesis atau kerangka konseptual masih parsial. & Sumber tidak relevan, tidak tersintesis, atau tidak berkontribusi pada kerangka konseptual. \\ \\
Kualitas rumusan masalah dan kerangka penelitian & 6 & Masalah dirumuskan secara operasional, terukur, dan kerangka penelitian konsisten dengan literatur yang ada. & Rumusan masalah cukup jelas tetapi kerangka penelitian masih lemah atau tidak konsisten. & Masalah tidak jelas atau tidak dapat diukur; kerangka penelitian tidak mendukung. \\ \\
Kesesuaian metode dan instrumen penelitian & 4 & Metode, instrumen, dan rencana validasi sesuai dengan jenis penelitian dan rumusan masalah yang dirumuskan. & Metode cukup sesuai tetapi instrumen atau prosedur validasi belum lengkap. & Metode tidak sesuai atau instrumen tidak dapat digunakan untuk menjawab masalah penelitian.\\}

\assessmentblock{Kuis Konsep Metodologi Penelitian}{Kuis tertulis}{SCPMK0403-03801, SCPMK0401-03804}{Mahasiswa mengerjakan kuis untuk mengukur pemahaman konsep jenis penelitian, etika, dan analisis data sesuai Sub-CPMK.}{Menjawab soal kuis secara mandiri dalam waktu yang ditentukan tanpa bantuan sumber terbuka.}{Lembar jawab kuis tertulis.}{Pemahaman jenis dan etika penelitian & 3 & Jenis penelitian, karakteristik, dan prinsip etika dijelaskan akurat dengan contoh relevan dari konteks TI. & Sebagian konsep benar tetapi penerapan atau contoh masih lemah. & Konsep keliru atau tidak dapat diterapkan pada konteks penelitian TI. \\ \\
Ketepatan identifikasi masalah dan hipotesis & 4 & Masalah diidentifikasi dari kasus, hipotesis dirumuskan secara operasional dan terukur sesuai kaidah ilmiah. & Identifikasi masalah cukup tepat tetapi hipotesis belum operasional atau tidak terukur. & Masalah tidak relevan atau hipotesis tidak dapat diuji secara ilmiah. \\ \\
Penerapan konsep analisis data & 3 & Jenis analisis (kuantitatif/kualitatif) dipilih tepat dan prosedurnya dijelaskan dengan benar sesuai data. & Pilihan analisis tepat tetapi prosedur atau justifikasinya belum lengkap. & Analisis tidak sesuai dengan jenis data atau masalah penelitian.\\}

\assessmentblock{UTS: Evaluasi Proposal Penelitian Awal (Bab 1 \& 2)}{Evaluasi portofolio proposal}{SCPMK0403-03801, SCPMK0403-03802}{Mahasiswa menyerahkan dan mempresentasikan draft proposal penelitian bab 1 dan 2 untuk dievaluasi kelengkapan dan kualitasnya.}{Menyusun draft proposal penelitian bab 1 (latar belakang, rumusan masalah, tujuan, manfaat) dan bab 2 (tinjauan pustaka, kerangka konseptual).}{Draft proposal bab 1 dan 2 dalam format dokumen ilmiah sesuai panduan JTI Polinema.}{Kelengkapan dan kualitas latar belakang serta rumusan masalah & 7 & Latar belakang menunjukkan masalah nyata berbasis data, rumusan masalah operasional, tujuan dan manfaat terukur. & Latar belakang dan rumusan masalah ada tetapi argumen atau data pendukung masih lemah. & Latar belakang tidak menunjukkan masalah nyata atau rumusan masalah tidak jelas dan tidak operasional. \\ \\
Kualitas tinjauan pustaka dan kerangka konseptual & 8 & Tinjauan pustaka dari sumber primer yang relevan, sintesis membentuk kerangka konseptual yang solid dan mendukung penelitian. & Tinjauan pustaka cukup tetapi sintesis atau kerangka konseptual belum kuat dan terhubung. & Tinjauan pustaka tidak tersistematisasi atau kerangka konseptual tidak mendukung penelitian. \\ \\
Konsistensi antara masalah, tujuan, dan rencana metode & 5 & Masalah, tujuan, pertanyaan penelitian, dan rencana metode saling terhubung dan konsisten satu sama lain. & Keterkaitan antara komponen ada tetapi belum sepenuhnya konsisten atau terdapat gap. & Komponen proposal tidak terhubung atau terdapat kontradiksi yang signifikan.\\}

\assessmentblock{PBL Proposal Penelitian Lengkap}{Project Based Learning}{SCPMK0403-03801--SCPMK0401-03804}{Mahasiswa menyusun dan mempresentasikan proposal penelitian lengkap (3 bab) sebagai luaran utama PBL semester.}{Menyusun proposal penelitian lengkap (bab 1--3), mempresentasikan di depan dosen dan/atau reviewer, serta merespons pertanyaan dengan argumentasi ilmiah.}{Dokumen proposal penelitian lengkap (bab 1--3) dan bahan presentasi sesuai format JTI Polinema.}{Kualitas proposal: sistematika, argumen, dan kontribusi ilmiah & 20 & Proposal memiliki sistematika lengkap, argumen kuat berbasis literatur, dan kontribusi ilmiah yang jelas dan terukur. & Proposal cukup sistematis tetapi argumen atau kontribusi ilmiahnya masih terbatas. & Proposal tidak sistematis, argumen lemah, atau kontribusi ilmiah tidak jelas. \\ \\
Ketepatan metodologi dan instrumen penelitian & 18 & Metodologi sesuai jenis penelitian, instrumen valid/reliabel, dan prosedur pengumpulan data terdokumentasi dengan baik. & Metodologi cukup sesuai tetapi instrumen atau prosedur pengumpulan data belum lengkap. & Metodologi tidak sesuai atau instrumen tidak dapat menjawab masalah penelitian yang dirumuskan. \\ \\
Kemampuan presentasi dan defensi proposal & 17 & Presentasi runtut, argumen dipertahankan dengan evidence, dan merespons pertanyaan dengan tepat dan percaya diri. & Presentasi cukup baik tetapi defensi atau respons terhadap pertanyaan masih lemah. & Presentasi tidak terstruktur atau tidak mampu mempertahankan argumen proposal penelitian.\\}

\assessmentblock{UAS Ujian Akhir Metodologi Penelitian}{Ujian tertulis}{SCPMK0403-03803, SCPMK0401-03804}{Mahasiswa mengerjakan ujian akhir semester untuk mengukur kompetensi akhir metodologi penelitian secara komprehensif.}{Menjawab soal ujian secara mandiri mencakup seluruh materi metodologi penelitian dalam waktu yang ditentukan.}{Lembar jawab ujian tertulis.}{Kemampuan analisis masalah dan desain penelitian & 8 & Masalah dianalisis secara komprehensif, desain penelitian dipilih dan diargumentasikan dengan tepat sesuai konteks. & Analisis masalah cukup tetapi pemilihan desain penelitian belum sepenuhnya diargumentasikan. & Analisis dangkal atau desain penelitian tidak sesuai dengan masalah yang diajukan. \\ \\
Penerapan analisis data dan interpretasi hasil & 7 & Metode analisis dipilih tepat, prosedur dijelaskan secara rinci, dan interpretasi hasil disajikan secara ilmiah. & Pilihan analisis tepat tetapi prosedur atau interpretasinya belum lengkap dan sistematis. & Analisis tidak sesuai atau interpretasi tidak berdasarkan data yang tersedia. \\ \\
Teknik penulisan karya ilmiah dan sitasi & 5 & Penulisan sesuai kaidah ilmiah, sitasi dan referensi mengikuti format yang ditetapkan secara konsisten. & Penulisan cukup ilmiah tetapi sitasi atau format referensi masih mengandung kesalahan. & Penulisan tidak mengikuti kaidah ilmiah atau sitasi tidak benar dan tidak konsisten.\\}

\begin{tblr}{width=\textwidth,colspec={|Q[l,m,3.2cm]|X[l,m]|},hlines,vlines,hspan=minimal,row{1-3}={bg=headgray},rowsep=1.5pt,colsep=3pt}
Jadwal Pelaksanaan: & Minggu 1--17 sesuai RPS. Setiap evaluasi memiliki RTM dan rubrik multi-kriteria. \\
Lain-Lain Yang Diperlukan: & LMS, perangkat reference manager, template proposal penelitian JTI Polinema, dan perangkat presentasi. \\
Pustaka: & Mengacu pustaka utama dan pendukung pada bagian RPS. \\
\end{tblr}
